\chapter{Cohen Up to Epsilon}

The previous chapters leaned, more and more openly, on a divide they never made
precise: countable exhaustion reaches the certificate family and stops short of the
coupling, and the two families were kept apart by type. That divide has the shape of a
set-theoretic fact, and a finite apparatus that means to approach the continuum should
put its set-theoretic house in order before it knocks on that door. This chapter does
the housekeeping. It builds the apparatus's own counting, files a finite forcing
calculus bounded by a tolerance, gives that calculus its order and its coding, and
closes the finite arc --- all before any continuum claim is made.

The chapter opens at the bottom, with counting. The apparatus has carried a number
type since its first pages, but it has never made that type count: never shown, in its
own number system, that one and one make two. The first section does exactly that. It
defines the successor and the sum in the apparatus's own numbers and proves the small
arithmetic that everything above will quietly use --- one and one is two, one and two is
three, one and one and one is three --- as definitional facts, with an induction
principle over the numbers the successor actually generates. It is the humblest section
in the book, and it is load-bearing: the closeout at the chapter's end will rest on a
single successor certificate proved here.

The second section introduces the objects the forcing calculus acts on. A spline
condition is a finite list of bounded knot-value assignments --- a finite piece of
partial information about where the spline sits, with every knot and value held below a
fixed tolerance. These are forcing conditions in the precise sense: finite, partial,
and orderable by how much they commit to. The section establishes that they behave: the
empty condition commits to nothing, and every condition has a finite support, the finite
set of knots it speaks about.

The third section gives the calculus its three operations, the three that any forcing
discipline needs. Extension is information growth: a stronger condition contains every
assignment of a weaker one, and extension is reflexive and transitive. Compatibility is
consistency: two conditions that agree wherever they overlap have a common extension,
their merge. And decision is density: for any bounded knot, a one-step extension decides
it, and any finite list of bounded knots can be decided at once. These three --- extend,
merge, decide --- are the Cohen-shaped content the apparatus can have before the continuum
door, and the section proves all of them stay finite, controlled by the same tolerance
the convergence argument used.

The fourth section gives the calculus its order and its size. The order is extension,
but conditions carry proof, so the order closes only up to the extensional equality of
the assignments they record; extension is a partial order up to that equality, and the
merge is its meet, commutative up to the same equality. And the calculus is countable:
each assignment carries a natural-number code, each condition carries a code built from
its assignments, and so the conditions are enumerable. This is the constructive content
of ``the ledger stays countable'' --- not a completed enumeration of a continuum, but a
finite code attached to each finite condition.

The fifth section closes the finite arc and hands the closeout to the gauge side. It
packages the forcing calculus, its order, and its coding into a single finite-compliance
certificate, and then it shows that the gauge residue rests on the counting of the first
section. From the one successor certificate --- that the successor of the successor of one
is the successor of two --- the apparatus reads off that the vacuum's linear term closes
and the residue is the mixed pair, the electron. The gauge commutator residue is closed
out against the count, and the finite arc that began with one and one making two ends
with the electron underwritten by that same arithmetic.

By the end of the chapter the apparatus has earned the right to approach the continuum.
Its counting is its own and proven; its forcing calculus is finite, ordered, countable,
and bounded by a tolerance; and the gauge residue is tied back to the count. Nothing in
the chapter crosses the continuum door. Everything in it makes sure that when the next
chapter does name what lies beyond that door, the house behind it is in order.

\section{Counting to Three in the Original Number System}

Before the continuum door opens, the apparatus has to put its finite ledger in
order. The previous chapter ended with a boundary reading that could not be
absorbed into an interior mixture. This chapter asks a different housekeeping
question: when the apparatus files finite conditions, extends them, checks
compatibility, decides bounded knots, and closes the gauge residue against a
count, what number system is doing the counting? The answer cannot be an
unexamined import from outside the apparatus. The ledger has to count in the
apparatus's own numbers.

That is why the section begins so low. The apparatus has carried a number type
since its earliest pages, but carrying a type is not yet having arithmetic. A
finite forcing condition will later be bounded by a count; a finite list of
knots will be decided by a count; the gauge closeout will rest on a small
successor certificate. All of that would be theatrical if the apparatus had not
first shown that its own numbers can reach the first nontrivial cases. Counting
to three is not cute. It is the floor under the finite Cohen-shaped ledger.

The construction starts with successor. There is a base number, written \(0\),
and there is an operation \(\mathrm{succ}\) that takes a number to its next
successor. The numerals used in the chapter are the successor-generated ones:
\[
  1=\mathrm{succ}\,0,
  \qquad
  2=\mathrm{succ}\,1,
  \qquad
  3=\mathrm{succ}\,2.
\]
Addition is then defined by recursion on the left argument:
\[
  0 + m = m,
  \qquad
  (\mathrm{succ}\,n) + m = \mathrm{succ}(n + m).
\]
This is the familiar recursive sum, but its importance here is local. The
operation is written inside the apparatus's number system, so the finite counts
that appear later are not smuggled in from a separate arithmetic background.

With successor and sum in place, the first identities are computations. The
apparatus does not need a long argument to see one and one become two; it
unfolds the definitions:
\[
  1 + 1 = 2,
  \qquad
  1 + 2 = 3,
  \qquad
  1 + 1 + 1 = 3.
\]
The equalities are definitional identities in the apparatus's own arithmetic.
They say that the three ways the chapter will name the first small counts agree:
two is \(1+1\), three is \(1+2\), and three is also \(1+1+1\). Later sections can
therefore speak about one-step and finite-step extensions without wondering
whether the underlying numerals have been aligned.

The point is not that three is a large number. It is that three is the first
place where the chapter's later bookkeeping begins to look like bookkeeping at
all. A condition can be empty, extended once, extended again, and compared with
another finite commitment; a small residue calculation can distinguish the
linear term, the closed vacuum part, and the mixed pair that remains. Those
uses do not require infinite arithmetic. They require the apparatus to know
which successor steps it has taken and to keep the names for those steps
coherent. The identities above give the reader that coherence before the
forcing language starts to spend it.

The next point is more delicate. The apparatus's number type is wider than the
successor-generated numerals. It may contain permissive values that are numbers
by type but are not reached by starting at \(0\) and applying successor. An
induction principle over the whole type would say too much. The honest
principle is induction along the canonical spine: the numbers generated from
the base by repeated successor. For any motive on numbers, the apparatus proves
the generated induction rule
\[
  \text{motive}(0),
  \quad
  \big(\forall n \text{ canonical},\ \text{motive}(n) \Rightarrow \text{motive}(\mathrm{succ}\,n)\big)
  \;\Longrightarrow\;
  \text{motive}(n) \text{ for canonical } n.
\]
The restriction matters. The finite ledger only needs the counts that the
apparatus can actually generate. Its later forcing conditions do not require a
global classification of every value admitted by the number type; they require a
reliable way to reason along the successor-built counts used in the finite
lists, finite supports, finite codes, and bounded decisions.

This also keeps the implication boundary clean. The apparatus is not claiming
that every permissive value in the number type has been enumerated by successor,
and it is not turning a finite count into a completed continuum listing. It is
separating the arithmetic it actually uses from the larger type that happens to
contain it. The later ledger may be countable and coded because each finite
condition receives a finite code; that is different from pretending that the
continuum has already been exhausted. The canonical spine is exactly the
right-sized domain for the finite counts that this chapter needs.

The section then isolates the successor certificate that the gauge closeout will
need. It is the small lifted identity
\[
  \mathrm{succ}(\mathrm{succ}\,1) = \mathrm{succ}\,2,
\]
obtained by applying successor to the identity that \(2=\mathrm{succ}\,1\). The
certificate is deliberately modest. It does not turn finite arithmetic into a
continuum enumeration, and it does not decide any set-theoretic question. It
only records that the apparatus can certify the exact successor step later used
when the gauge residue is closed against the count.

This is the section's discipline. Before the finite forcing calculus talks
about extension, compatibility, decision, order, coding, and residue closeout,
the apparatus establishes the tiny arithmetic those words will spend. The
numbers are successor-generated. Addition is recursive. The identities
\(1+1=2\), \(1+2=3\), and \(1+1+1=3\) compute in the system itself. Induction is
available on the canonical spine and no wider than that. The final successor
certificate is filed for later use. The chapter can now raise its finite
Cohen-shaped ledger on an arithmetic floor the apparatus has actually laid.

\section{Finite Spline Conditions}

The previous section put the counting floor in place. The apparatus can name
the first successor-generated numbers, add them in its own arithmetic, and
reason along the canonical spine. Now the chapter can introduce the finite
objects that will spend those counts. A forcing ledger does not begin with a
completed spline. It begins with a finite piece of information: this knot has
this value, that knot has that value, and nothing has yet been said about the
rest. This section gives that finite piece of information its name.

The name is \emph{spline condition}. The phrase is deliberately modest. A
spline condition is not the spline, not a total function, and not an analytic
object stretched across a continuum. It is a finite record of commitments. Each
commitment says that at one named knot the spline is assigned one named value.
The apparatus keeps only finitely many such commitments at a time, and every
one of them is held below the same finite allowance. That is the whole point:
the condition is small enough to be inspected, counted, extended, and compared
inside the apparatus before any continuum door opens.

Fix a finite tolerance, written \(\varepsilon\). In the chapter's internal
language this tolerance is a number in the apparatus's finite scale. It is not a
metric completion and not a hidden real-valued limit. It is the same sort of
finite allowance that Chapter 14 used when the late gauge readings were forced
into a shrinking margin. Here the allowance has a bookkeeping role. It tells the
forcing ledger which knots and values are small enough to be admitted into a
condition.

The smallest commitment is an atom. An atom is a knot-value assignment:
\[
  \text{atom} = (\mathrm{knot}, \mathrm{value}),
  \qquad
  \text{bounded by } \varepsilon
  \;\Longleftrightarrow\;
  \mathrm{knot} \le \varepsilon \ \text{ and }\ \mathrm{value} \le \varepsilon.
\]
The atom therefore carries two pieces of information: which knot is being
mentioned, and which value is being assigned there. Boundedness is not an
external comment added after the fact. It is part of what makes the atom
admissible for this finite calculus. If the knot lies beyond the allowance, or
if the value lies beyond the allowance, the atom is not one of the bounded atoms
this section is allowed to file.

The bounded atom is the first place where the chapter's two disciplines meet.
The knot side says where the finite record speaks. The value side says what the
record asserts there. The tolerance says that both coordinates are still inside
the finite range the apparatus has agreed to handle. None of this asks for a
dense line of knots or for a function defined at every possible location. The
atom is a single local commitment. It can be written down, checked, and counted.
That is why it is the right unit for a finite forcing ledger.

The chapter also keeps the atom separate from interpretation. A bounded atom
does not yet say whether the chosen value is physically realized, whether the
condition will survive later comparison, or whether a future extension will add
neighboring knots. It only says that this finite row is admissible as a row. The
later calculus may extend, compare, and decide such rows, but those operations
need an object simple enough to receive them. The atom supplies that object: one
knot, one value, one boundedness check.

A spline condition is a finite list of bounded atoms:
\[
  C = [\,\text{atom}_1, \ldots, \text{atom}_k\,],
  \qquad
  \text{every atom bounded by } \varepsilon.
\]
The word ``list'' is important. The apparatus is not quantifying over all
knots; it is holding a finite record. The length of the list is a count the
apparatus can understand. The atoms in the list are the commitments currently
made. The boundedness proof travels with the record, so a condition is not
merely a list that happens to be small in the reader's imagination. It is a list
equipped with the guarantee that each of its atoms respects \(\varepsilon\).
Later sections can therefore manipulate conditions without rechecking from
scratch whether the data have wandered outside the finite allowance.

A condition may mention one knot, many knots, or no knots at all, but always
only finitely many. Its list structure is the reason the chapter can later talk
about finite support and natural coding. A page with finitely many rows can be
copied, extended by adding another row, inspected for repeated knots, and
assigned a finite code. A total spline cannot be handled in that way inside the
present chapter. The spline condition is therefore not a poor substitute for the
total object; it is the finite object the apparatus actually knows how to use.

The condition also keeps its boundedness pointwise. The apparatus does not
declare a whole condition bounded by intuition or by a global promise. It
checks the atoms one by one. If the first atom is bounded, and the second atom
is bounded, and so on through the finite list, then the condition is bounded.
This pointwise structure is what makes the finite list stable under later
bookkeeping. Whenever a later step adds one more bounded atom, the old finite
bounded record can be enlarged without changing what it already certified.

The empty list is already a condition:
\[
  \varnothing = [\,],
  \qquad
  \text{bounded by } \varepsilon \text{ vacuously.}
\]
It commits to no knot and assigns no value. That is not a defect. In a forcing
calculus the empty condition is the least initial commitment, the record before
any finite information has been added. Its boundedness is vacuous because there
are no atoms that could violate the tolerance. This gives the ledger a clean
starting point: one may begin with no commitments at all and then add finite
information by controlled steps.

Every nonempty condition speaks only where its atoms speak. Its finite support
is the list of knots mentioned by its atoms:
\[
  \mathrm{support}(C) = [\,\mathrm{knot}(\text{atom}) : \text{atom} \in C\,],
  \qquad
  \text{finite.}
\]
Because \(C\) is a finite list, its support is finite as well. Repetitions and
proof-carrying details can be handled later by the order and coding sections;
for the present section the essential fact is simpler. A condition has a finite
domain of speech. It constrains the knots in its support and remains silent
elsewhere.

The support forgets the values and remembers only where the condition has
spoken. This is a useful loss of information. If one wants to know which knots
remain open, the support is the list to consult. If one wants to know whether
two conditions might clash, the first question is whether their supports
overlap. If one wants to code a condition later, the support is part of the
finite data that makes coding possible. The support is not an extra continuum
domain hiding behind the condition. It is just the finite list of mentioned
knots, extracted from the finite list of atoms.

The support also explains why partiality is precise rather than vague. A
condition is not merely "incomplete" in a casual sense. It has an exact finite
frontier: the knots in the support have been mentioned, and the knots outside
the support have not. This lets the apparatus distinguish silence from
disagreement. A missing atom at a knot is not a contrary assignment; it is no
assignment. That distinction will matter when the next section asks whether two
finite records can stand together.

That silence is load-bearing. Outside the support, the condition says nothing.
It does not secretly assign a default value, and it does not pretend to know the
whole spline. The unmentioned knots are still open. This is the honest
partiality that makes forcing possible: a condition can be extended precisely
because it has not already answered every question. If a later step wants to
decide a new bounded knot, it can add a bounded atom at that knot. If a later
step wants to compare two pieces of information, it can ask whether their
finite commitments clash where their supports overlap. The present section
does not perform those operations, but it gives them the right kind of object.

The tolerance discipline matters just as much as finite support. A finite list
could still be useless for this chapter if it contained an unbounded knot or an
unbounded value. Such a list would no longer belong to the controlled ledger.
By requiring every atom to sit below \(\varepsilon\), the apparatus prevents
the finite calculus from drifting away from the same scale that governed the
late convergence discussion. The connection to Chapter 14 is only a discipline,
not a reopened proof. The earlier chapter used a vanishing geometric tolerance
to squeeze readings together. This chapter fixes a finite allowance so that its
partial records remain controlled while the set-theoretic bookkeeping is done.

The section can now say what a condition is and what it is not. It is finite:
the apparatus can count its atoms, inspect them, and collect their knots. It is
bounded: every atom stays under \(\varepsilon\). It is partial: outside its
support it makes no claim. It is data with a certificate of boundedness, not a
guess about a total spline. These four features are exactly what the next
section will need when it defines extension, compatibility, and decision.
Extension will mean adding information without losing the old commitments.
Compatibility will mean that two finite records can stand together when they
do not disagree on their overlap. Decision will mean that a bounded knot can be
answered by a finite controlled extension. Those are S03's operations; this
section keeps to the underlying objects.

It is useful to picture a condition as a ledger page with finitely many rows.
Each row has a knot, a value, and the evidence that both lie within the finite
allowance. Blank rows are not hidden assertions. They are places where the
ledger has not yet spoken. A longer ledger page may later include more rows,
and two ledger pages may later be compared or merged if their rows do not
contradict one another. But at this stage the page itself is already a valid
mathematical object. It is not waiting for the completed spline to make it
meaningful.

This is the finite Cohen-shaped posture of the chapter. The apparatus treats
partial information as real information. A condition's incompleteness is not
embarrassing; it is the feature that leaves room for extension. Its boundedness
keeps the ledger below the chosen allowance. Its finite support keeps the
record inspectable. Its silence outside the support keeps the continuum door
closed. Before the next section lets conditions grow, combine, and decide, this
section has made sure that the things being manipulated are finite, bounded,
partial, and honest about what they do not say.

This is also why no continuum claim has entered. The section has not produced a
generic spline, enumerated all possible spline values, or shown that the finite
ledger exhausts an analytic space. It has only prepared the finite records that
the apparatus can responsibly manipulate. Each record is a bounded list of
bounded atoms with finite support. That is enough for the next section's
forcing-shaped moves and no more than enough. The chapter stays inside its
promise: finite set-theoretic hygiene up to \(\varepsilon\), with the continuum
door still closed.

\section{Extension, Compatibility, and Decision}

The previous section introduced the objects of the ledger: finite lists of
bounded atoms under one fixed finite allowance. This section gives those
objects their finite motion. A condition can grow by adding information. Two
conditions can be compared for consistency on the knots where they both speak.
Compatible conditions can be merged into a common stronger condition. And a
bounded knot can be decided by a finite extension. These are the operations that
make the ledger Cohen-shaped, but every one of them stays inside the finite
apparatus.

The first operation is extension. Extension is information growth without
forgetting. A condition \(C'\) extends a condition \(C\) when every commitment
made by \(C\) is still present in \(C'\):
\[
  C' \text{ extends } C
  \quad\Longleftrightarrow\quad
  \text{every atom of } C \text{ is an atom of } C'.
\]
The stronger condition may add more bounded atoms, but it may not retract the
old ones. In ledger language, \(C'\) contains all rows already written in \(C\)
and perhaps some additional rows. The word "stronger" therefore means "more
committed", not more mysterious.

Two elementary facts follow from the definition. Extension is reflexive: every
condition extends itself, because every atom of \(C\) is already an atom of
\(C\). Extension is transitive: if \(C''\) contains every atom of \(C'\), and
\(C'\) contains every atom of \(C\), then \(C''\) contains every atom of \(C\).
Thus the ledger has the basic direction required of a forcing order. Moving
down to a stronger condition means adding finite bounded commitments while
preserving the finite bounded commitments already made.

The preservation clause is what keeps extension honest. A new condition cannot
decide a knot by changing what the old condition said there. It can only keep
the old assignment or add assignments where the old condition was silent. If a
later proof carries proof-laden boundedness data along with the list, that data
travels with the extension. The section does not need the full extensional
order yet; S04 will handle proof-carrying sameness and coding. Here the
reader-facing point is the finite one: extension is controlled information
growth.

This also explains why extension is the right direction for a forcing ledger.
The weaker condition is the less informative one. It leaves more knots open.
The stronger condition is the one with more finite commitments. Nothing about
the word "stronger" means larger in value, closer to a limit, or more
continuum-like. Strength is commitment. A page with three bounded rows is
stronger than the same page with two of those rows, because it has answered one
more finite question while keeping the earlier answers. That is the exact
information order the later finite calculus will use.

The empty condition illustrates the direction. Since the empty condition has no
atoms, every condition extends it. The empty page is weakest because it says
nothing. A one-row condition is stronger because it has made one bounded
commitment. A two-row condition extending that one is stronger again. This
chain is still finite at every step. The apparatus can follow it without
pretending that there is already an infinite completed record waiting at the
end.

The second operation is compatibility. Compatibility is consistency on overlap.
Two conditions \(C\) and \(D\) are compatible when they do not assign different
values to the same knot. Equivalently, whenever a knot lies in both finite
supports, the two conditions give it the same value. A shared knot with matching
values is harmless. A shared knot with conflicting values blocks combination.
There is no finite condition that can honestly contain both assignments.

When two conditions are compatible, their merge is the finite record obtained
by putting their assignments together. In symbols:
\[
  C, D \text{ compatible}
  \;\Longrightarrow\;
  C \cup D \text{ extends } C \ \text{ and }\ C \cup D \text{ extends } D.
\]
The merge is finite because it is made from two finite lists. It is bounded
because every atom in either list was already bounded by the same
\(\varepsilon\). It extends \(C\) because all rows of \(C\) remain present; it
extends \(D\) for the same reason. The only thing compatibility has to prevent
is contradiction at a shared knot. Once that contradiction is absent, the
finite records can stand together.

The merge should be understood as a common extension, not as a new kind of
object. It is still a spline condition: a finite list of bounded atoms. It is
common because it extends each input condition. It is an extension because it
keeps every old assignment. It is bounded because both inputs used the same
\(\varepsilon\), so no atom introduced by the merge lies outside the finite
allowance. The operation does not average values, interpolate between knots, or
invent a total spline. It simply places two consistent finite records on one
page.

This common-extension view is the finite substitute for the more dramatic
language that forcing can invite. The apparatus is not saying that compatible
conditions reveal a hidden global object. It is saying that, when two finite
pieces of information do not contradict each other, there is a finite piece of
information that contains them both. That is enough for the local calculus. The
later order section may describe this merge as a meet up to extensional
equality, but the present section needs only the concrete fact: compatible
finite records can be jointly preserved by another finite record.

The converse is just as important. If \(C\) and \(D\) disagree at a knot in
their overlap, then no common extension can exist. A common extension would have
to preserve the assignment from \(C\) and also preserve the different assignment
from \(D\), thereby assigning two values to one knot. That is not a stronger
condition; it is an inconsistent record. Thus compatibility is exactly the
finite test for whether two partial records can have a common stronger record.
Consistent partial information combines; inconsistent partial information does
not.

The third operation is decision. A condition decides a bounded knot \(k\) when
it contains some bounded atom assigning a value at \(k\). It does not decide
all knots, and it does not complete the spline. It answers one finite question:
what value has this condition committed to at this bounded knot? If \(k\) is
already in the support, the condition has already spoken there. If \(k\) is
outside the support, the condition is silent there and can be extended by adding
one new bounded atom.

The one-step decision move is the finite density fact. Given a condition \(C\)
and a bounded knot \(k\), choose a bounded value allowed by the tolerance and
add the atom \((k,\mathrm{value})\) when it does not conflict with what \(C\)
already says. The result is a finite condition \(C'\) extending \(C\), and
\(C'\) decides \(k\):
\[
  \text{for any condition } C \text{ and bounded knot } k,
  \quad
  \exists\, C' \text{ extending } C,\quad C' \text{ decides } k.
\]
If \(C\) already assigns a value at \(k\), the decision is already present. If
the new choice would contradict an old assignment, choose the old assignment
instead. In either case, the decision remains a finite bounded commitment. No
infinite object is summoned to answer a one-knot question.

The decision move is intentionally local. It does not ask the apparatus to know
which value the final spline "really" has at every knot. It asks only for a
bounded commitment at the particular bounded knot under discussion. If the
ledger has not yet spoken there, the apparatus may add one bounded row. If it
has already spoken there, the old row is the decision. In both cases the answer
is found inside a finite extension of the current condition. This is why the
operation deserves the forcing word "dense" while remaining finite: below the
current finite record, there is a stronger finite record that answers the
particular finite question.

The same idea iterates across a finite list of bounded knots. Decide the first
knot by one finite extension. Decide the second by extending that result. Keep
going until the finite list is exhausted. Because the list is finite, the
process takes finitely many steps. Because each new atom is bounded by
\(\varepsilon\), the final condition remains bounded. Because each step
extends the previous condition, no earlier commitment is lost. The apparatus
can therefore settle any finite list of bounded questions by finitely many
bounded extensions.

The finite-list qualification is the guardrail. If the apparatus is handed
three bounded knots, it can make at most three one-step additions, skipping any
knot already decided. If it is handed ten bounded knots, the same reasoning
takes ten finite steps. The count may vary, but it remains a count the
apparatus can carry. What the section does not claim is a single completed
condition that has answered every possible bounded knot in advance. The
apparatus answers the questions actually placed on the finite ledger.

This is the density available to the finite apparatus. It is strong enough for
the chapter's ledger: every actually asked finite question can be answered
below some stronger finite condition. It is not the full generic object of
set-theoretic forcing. It does not decide every possible knot at once, and it
does not create a total spline. It says only that finite questions, asked
finitely, can be answered by finite controlled growth. That is the exact level
of forcing behavior the chapter is allowed to claim before the continuum door.

The three operations also preserve the two controls from S02. Extension
preserves finiteness because it adds only finitely many atoms in the situations
used here; it preserves boundedness because the added atoms are bounded.
Compatibility checks only a finite overlap of finite supports. Merge combines
two finite bounded records into another finite bounded record. Decision adds
one bounded row, or finitely many bounded rows for a finite list of questions.
At no point does the calculus escape into an unbounded assignment or a
decided-everywhere spline.

It is useful to keep the ledger picture in view. Extension writes more rows
without erasing the old ones. Compatibility asks whether two pages disagree on
any row whose knot appears on both pages. Merge copies two nonconflicting pages
onto a common page. Decision adds the row needed to answer a particular bounded
question. These are ordinary finite bookkeeping moves, but together they give
the chapter its Cohen-shaped discipline: partial information can grow, combine,
and settle questions while remaining partial.

The restraint is part of the result. A full set-theoretic forcing argument
eventually speaks about a generic object that meets every dense requirement and
settles an infinite field of questions. This chapter does not take that step.
It keeps the finite conditions and the finite operations. The generic object is
named only as the thing beyond the door, not as something already constructed
inside the apparatus. The continuum is treated the same way. The finite ledger
can extend, merge, and decide below \(\varepsilon\), but it does not become a
completed continuum enumeration.

Thus S03 supplies the active content of the finite forcing calculus. S02 gave
the objects: bounded finite partial conditions with finite support. This
section gives the moves: preserve old commitments by extension, combine
compatible records by merge, and decide bounded finite questions by bounded
finite extension. S04 will put these moves into an extensional order and attach
natural codes. S05 will return to the gauge residue closeout. For now the
chapter has exactly what it needs: a finite calculus in which partial
information grows honestly and stops before it pretends to be total.

\section{Extensional Order and Natural Coding}

S03 supplied the finite operations: extension, compatibility, merge, and
decision. This section packages those operations as an order and a code. The
order says when one finite record is stronger than another. The code says how
the family of finite records can be listed by natural numbers. Together they
give the chapter the promised countable forcing shape, but still only at the
level of finite bounded conditions. Nothing here crosses into a completed
generic spline or an enumeration of continuum values.

The first subtlety is equality. A condition is not merely a bare set of pairs in
the apparatus. It may carry boundedness certificates, support witnesses, or
other finite evidence showing that its atoms really sit below the fixed
\(\varepsilon\)-allowance. Two records can therefore have different internal
packaging while making exactly the same commitments. For the forcing order, the
right equality is extensional equality:
\[
  C \equiv D
  \quad\Longleftrightarrow\quad
  \text{\(C\) and \(D\) record the same bounded atoms.}
\]
The condition is read by what it commits to at bounded knots, not by the route
by which the apparatus certified those commitments.

This is not a relaxation of the finite discipline. It is the discipline stated
at the correct level. If \(C\) and \(D\) carry the same atom \((k,v)\) for
every knot-value commitment, then no later finite extension, compatibility
test, or decision test can distinguish them by their commitments. They are the
same partial information. Their certificates may have been assembled in
different ways, but those certificates support one finite record. The forcing
calculus must therefore order conditions extensionally: by preserved
assignments.

With that equality understood, extension becomes a partial order. Reflexivity
is immediate: every condition preserves all of its own atoms. Transitivity is
the same preservation read through two steps: if \(C''\) preserves all atoms of
\(C'\), and \(C'\) preserves all atoms of \(C\), then \(C''\) preserves all
atoms of \(C\). Antisymmetry holds extensionally. If \(C\) extends \(D\) and
\(D\) extends \(C\), then each condition contains every atom of the other. They
record the same commitments, so \(C \equiv D\).

The order is therefore a genuine forcing order once equality is read
extensionally. The weaker record is lower in information. The stronger record
has all the weaker commitments and perhaps more. A condition is not stronger
because it is closer to a continuum object or because its values are larger. It
is stronger because it has answered more bounded finite questions without
erasing earlier answers. This is exactly the direction S03 used for decision:
to decide a bounded knot, move to a stronger finite record.

The phrase "up to extensional equality" should be read as a bookkeeping
precision, not as a loophole. The apparatus may need several finite witnesses
to know that a row is allowed: a witness that the knot belongs to the finite
support, a witness that the value lies inside the current tolerance, or a
witness that adding the row preserves compatibility. Those witnesses can differ
while the row they certify is the same row. Antisymmetry cannot reasonably ask
the witnesses to be literally the same object. It asks the commitments to be
the same. If each record extends the other, every committed atom on either side
is present on both sides, so the two records have the same observable finite
content. That is the order-theoretic fact the forcing calculus needs.

This also keeps the notation honest. The chapter can write \(C' \leq C\), or
speak of \(C'\) as stronger than \(C\), once the convention is fixed: stronger
means extension by preserved finite information. If two presentations differ
only by certificate packaging, they occupy the same position in the order.
Future uses of the order should therefore never depend on the accidental route
by which the apparatus assembled a condition. They may depend only on which
bounded atoms the condition records and which atoms a stronger condition
preserves.

Compatibility and merge also become order statements. If two conditions agree
on overlap, their union is a common stronger condition. It extends the first
because it preserves every atom of the first, and it extends the second for the
same reason. It is least among common stronger conditions up to extensional
equality: any condition that extends both inputs must contain every atom from
both, and therefore must contain the merged record. The merge is not an
interpolation and not a value average. It is the least finite page that keeps
both nonconflicting pages.

This least-common-extension claim is deliberately modest. It applies when the
inputs are compatible finite records under the same finite allowance. If the
records disagree at a shared knot, there is no common stronger condition,
because a stronger condition would have to preserve both incompatible
assignments. If the records agree, the merge is finite because it combines
finite lists, bounded because every atom already obeyed the shared
\(\varepsilon\), and least because any common extension must include those
same atoms. The order does no hidden analytic work.

Extensional equality also keeps merge from depending on list order or
certificate order. Writing the atoms from \(C\) before the atoms from \(D\), or
from \(D\) before the atoms from \(C\), produces the same commitments. Removing
duplicate matching atoms changes no commitment. Reusing a different proof that
the same atom is bounded changes no commitment. All of these presentations are
one condition for order purposes. The apparatus is allowed to remember its
certificates internally while the forcing calculus reads only the finite
commitments those certificates certify.

The second structure is a natural-number code. A countable forcing needs its
conditions to be listable. Here listability is not asserted abstractly; it is
built from the finite data already present. A bounded atom has a knot and a
bounded value. The knot data can be assigned a natural code. The value data,
because it is admitted only through the finite bounded apparatus of the
chapter, can also be assigned a natural code. Pair the two codes to obtain a
code for the atom:
\[
  \mathrm{code}(k,v)
  =
  \langle \mathrm{code}(k), \mathrm{code}(v) \rangle_{\mathbb N}.
\]
The angle brackets here mean any fixed finite pairing operation on natural
numbers. The particular pairing convention is not the content; the content is
that a pair of natural codes receives one natural code.

A condition is a finite list of such atoms. Finite lists can be coded
recursively: the empty list receives a base code, and a nonempty list receives
a code built from the code of its head atom and the code of its tail list. Thus
each finite condition receives a natural number:
\[
  \mathrm{code}(\emptyset) \in \mathbb N,
  \qquad
  \mathrm{code}((a_0,\ldots,a_{n-1})) \in \mathbb N.
\]
If the apparatus uses a canonical ordering of the finite support, the code can
be made insensitive to the accidental order in which atoms were written. If it
does not canonicalize first, the resulting code still lists presentations, and
extensional equality then identifies presentations with the same commitments.
Either route stays finite.

One can picture the coding as a ledger index. First number the admissible knot
descriptions. Then number the admissible bounded value descriptions. Then use a
pairing operation to number each atom. Finally, use a finite-list code to number
the page of atoms. The construction is primitive in spirit: pair, cons, stop.
It does not require the apparatus to inspect infinitely many possible rows
before coding the row in hand. Given a finite condition, the apparatus reads
the finitely many atoms actually present and produces one natural number for
that finite page.

If duplicate or permuted presentations are allowed, the code may count more
presentations than extensional conditions. That is harmless for countability. A
surjection from natural-number codes onto presentations already gives a
listing, and quotienting presentations by extensional equality only identifies
some entries of that list. If the apparatus wants an injection instead, it can
choose a canonical presentation: sort the finite support by its knot code,
discard duplicate matching atoms, and reject any conflicting duplicate as an
invalid condition. Either method yields the same mathematical boundary: the
finite condition family is countable because finite coded data can be listed.

The countability conclusion is now concrete. Conditions inject into, or are
listed by, natural-number codes once the chosen presentation convention is
fixed. There is no appeal to an external continuum, no appeal to a completed
real line, and no appeal to an unbounded search through possible splines. The
ledger has only finite records, and finite records can be encoded by natural
numbers. That is the constructive content of saying that this forcing poset is
countable.

The qualification matters. Countability belongs to the family of finite
conditions, not to a completed generic object. A forcing argument may later
speak as though a generic object meets every dense requirement, but this
chapter has not built such an object. It has built a countable supply of finite
bounded approximants. The fact that these approximants can be listed does not
mean the continuum has been listed. It means the apparatus knows how to index
the finite questions and finite answers it is actually allowed to manipulate.

Nor does natural coding turn the values into ordinary real numbers. The code is
attached to the bounded finite data admitted by the chapter: knots, values, and
certificates inside the tolerance regime. Coding an admitted finite value is
bookkeeping. It does not complete a metric space, add limits, or certify that
all analytic values have been exhausted. The code says, "this finite atom has a
place in the ledger." It does not say, "the continuum has been captured."

The order and the code also fit each other. Extension can be checked against
the finite atom lists named by the codes. Compatibility can be checked by
looking through the finitely many coded knots in each support. Merge can be
recoded as the finite list obtained by combining compatible coded atoms.
Decision adds a new coded atom, or recognizes that the code for the required
atom is already present. Every operation from S03 therefore has a
natural-number shadow, but the shadow follows the finite operation rather than
replacing it.

That last direction matters. The code is not the condition's source of truth.
The finite commitments are the source of truth; the code is their index. A bad
code, a noncanonical code, or two codes for extensionally equal presentations
does not change which knots have been committed. The order is defined on the
commitments and then represented by codes. This prevents the natural-number
apparatus from smuggling in a stronger claim than the finite ledger earned.
Coding is a way to list and revisit finite commitments, not a way to enlarge
them.

This is why the chapter can call the structure Cohen-shaped without
overclaiming. It has a partially ordered family of finite conditions. It has
finite extension, finite compatibility, finite merge, and finite decision. It
has a natural coding that lists the finite conditions. These are the parts of
Cohen forcing the apparatus can responsibly instantiate below
\(\varepsilon\). The parts that require an actual generic object remain beyond
the finite door.

The completed package is a countable poset of finite bounded commitments,
ordered by extension up to extensional equality. Compatible records have least
common stronger records, again up to what they actually commit to. Each atom
and each condition has a natural code. The family of conditions is therefore
countable in the constructive sense available to the apparatus. What has been
counted is not the continuum, but the finite ledger of possible bounded
commitments.

S05 can now return to the gauge side. The present section has finished the
set-theoretic hygiene: conditions, operations, order, and code. The calculus is
ordered enough to behave like forcing and small enough to remain a finite
apparatus. The residue closeout that follows should not reopen the order or
the coding. It can use the package as established: a countable finite forcing
calculus, bounded by \(\varepsilon\), honest about its partiality, and stopped
before it mistakes its finite codes for a continuum.

\section{The Gauge Commutator Residue Closeout}

The numbered chapter closes by refusing to let its two halves drift apart. The
first half built the local arithmetic floor: successors can be counted inside
the apparatus, finite lists can be held as finite lists, and bounded questions
can be asked without pretending that a completed continuum has arrived. The
second half built the forcing-facing package below \(\varepsilon\): finite
conditions, finite extension, compatibility, merge, finite decision,
extensional order, and natural-number codes. S05 uses that package once, and
only once, to return to the gauge commutator residue. It is a closeout, not a
new theorem.

The certificate assembled by S04 has three faces. First, it is finite: every
condition is a finite record of bounded atom commitments. Second, it is ordered:
extension means adding information while preserving the older commitments, and
the order is read up to extensional equality. Third, it is countable in the
constructive sense available here: atoms have natural codes and finite
conditions have codes built from finite lists of atoms. Taken together, these
faces give the chapter its Cohen-compliance certificate below the tolerance. The
certificate says that the apparatus has a finite forcing calculus with the
right order behavior and the right coding behavior. It does not say that a
generic spline has been formed, that the continuum has been enumerated, or that
any completion hypothesis has been discharged.

This restraint is exactly what lets the gauge closeout be honest. The gauge
side does not need a continuum object from this chapter. It needs a finite
bookkeeping result: after the linear vacuum certificate has closed, the
surviving commutator contribution is the mixed-pair residue. In the language of
the earlier boundary chapters, this is the electron residue. The present chapter
does not rederive the analytic boundary story. It supplies the finite counting
and finite-condition hygiene that keep the residue from floating free of the
apparatus that names it.

The counting floor is the hinge. The apparatus first had to certify that its
successor arithmetic could reach the small finite counts it uses: one, its
successor, and the next successor. That is the humble certificate behind the
later bookkeeping. Once the finite count is available, the apparatus can sort
the terms it has already licensed: the vacuum's linear part closes under the
finite certificate, while the mixed pair is not absorbed into that linear
closure. What remains is the residue. Schematically, and only as a finite
accounting statement, the closeout has the form
\[
  \text{successor/counting certificate}
  \;\Longrightarrow\;
  \text{linear vacuum closure}
  \quad\text{with}\quad
  \text{residue} = \text{mixed pair}.
\]
The displayed implication is not an analytic completion theorem. It is the
chapter's bookkeeping join: the finite count supports the linear closeout, and
the unabsorbed mixed pair is the finite residue carried forward from the gauge
analysis.

The extensional order and the natural codes matter here because they prevent a
subtle overstatement. A residue cannot be certified by an amorphous pile of
finite records. The apparatus must know when two records make the same bounded
commitments, when one record extends another, when two records can be merged,
and how the resulting finite record is counted. Those are precisely the S03 and
S04 operations. They let the finite certificate be filed without treating its
route, proof packaging, or later use as part of the object certified. The record
that supports the gauge closeout is the finite record of commitments itself.

Thus the residue closeout is conditional on the apparatus already built in this
chapter. Given the successor/counting certificate, given finite bounded
conditions below \(\varepsilon\), given the extension and compatibility
calculus, and given the extensional coding of finite conditions, the chapter can
file the finite closure claim: the linear/vacuum part has closed in the finite
ledger, and the mixed-pair contribution is the remaining gauge residue. That is
the most this chapter claims. It names no new Sobolev space, asserts no Hilbert
completion, and supplies no completed generic object.

The chapter therefore ends with a deliberately finite package. On the
set-theoretic side, it has a countable family of finite Cohen-style conditions,
ordered extensionally and coded by natural numbers. On the gauge side, it has a
finite residue statement tied back to the counting floor. The two are joined by
bookkeeping rather than by a new continuum principle: finite records can be
counted, finite extensions can close finite questions, and the finite gauge
ledger can say what remains after its linear part is closed.

There is also a useful negative statement in the closeout. No part of the
argument asks the coding to choose a generic filter, no part asks the finite
order to supply a total assignment, and no part asks the gauge residue to carry
more analytic structure than the earlier chapters assigned to it. The residue
is filed after the finite certificates have done their finite work.

This is why the bridge may now approach the completion door without smuggling
anything through it. The numbered chapter has put the finite house in order. It
has not crossed to a completed continuum, a total generic spline, or a completed
boundary-radiation theorem. It has prepared the finite certificate that any
later completion discussion must cite: counting at the floor, Cohen-style
finite compliance in the middle, and the mixed-pair gauge residue at the top.
The next section may name the completion hypotheses because this one has
stopped exactly where the finite apparatus stops.

\section*{Bridge: Approaching the Completion Door}

The chapter has now done the finite work it promised to do. It began by making
the apparatus count in its own numbers. It then turned that count into a ledger
of finite spline conditions: bounded atoms, finite records, and partial
commitments below \(\varepsilon\). It showed how one condition extends another,
how compatibility is tested on overlap, how compatible records merge, and how a
finite list of bounded questions can be decided by finite extension. It then
quotiented away certificate packaging by extensional equality and coded the
finite records by natural numbers. Finally, it returned that finite package to
the gauge side and filed the mixed-pair residue after the linear/vacuum part had
closed.

That is a substantial package, but it is deliberately not a completion. The
chapter has a counting floor, a finite condition calculus, an extensional order,
a natural coding, and a finite gauge residue closeout. It does not have a total
generic spline. It does not have a completed enumeration of continuum objects.
It does not have a Hilbert space completion, a Sobolev completion, or a
continuum-facing result. The distinction matters because the book has now reached the
point where a careless sentence could make the finite apparatus sound stronger
than it is. Chapter 15's strength is exactly its refusal to do that.

The door is nevertheless visible. The finite conditions behave like the finite
traces of something larger. Their natural codes give a countable ledger, but a
countable ledger of finite records is not the same as a completed object. Their
extension order lets the apparatus answer any actually asked finite question,
but answering finite questions one by one is not the same as producing a total
assignment. Their residue closeout ties the gauge charge back to the count, but
that bookkeeping closeout is not the same as an analytic boundary-radiation
theorem. The bridge must say all of this plainly because the next chapter can
only be honest if the threshold is honest.

The threshold has appeared before under several names. The square-root door of
the operator discussion, the completion slot of the Hilbert frame, the generic
object of the forcing discussion, and the continuum-facing boundary family are
not four unrelated mysteries. They are different views of the same gap between a
finite apparatus and the completed structure it points toward. Chapter 15 has
not closed that gap. It has prepared the finite side of it carefully enough that
the next chapter can name what remains without handwaving.

The preparation is also negative preparation. By the end of the numbered
sections, the apparatus knows exactly which moves it is not allowed to treat as
free. A finite code is not a choice principle. A sequence of finite extensions
is not a completed filter. A residue statement certified by finite bookkeeping
is not a boundary analysis in a completed normed space. These refusals are not
timidity; they are what make the next chapter legible. Once the finite side has
said no to these substitutions, the completion side can ask for them openly and
measure what each one would add.

Chapter 16 therefore has a narrower job than triumph would suggest. It must
name the Sobolev/Hilbert door as a door, not as something already crossed. It
must identify the separable pre-Hilbert skeleton already suggested by the finite
energy bookkeeping, while distinguishing that skeleton from a completed Hilbert
space. It must explain how finite certificate families shadow continuum
certificate families without pretending the shadow is the object. It must bring
orthogonality, Clifford compatibility, and the boundary trace into one
conditional frame. And it must state generalized boundary radiation in its true
form: a theorem only under the completion hypotheses that make its analytic
words meaningful.

This bridge does not prove those obligations. It only aligns them. The finite
apparatus has earned the right to ask for completion hypotheses because it has
kept its own finite account clean. The next chapter will ask which hypotheses
are needed, which finite structures they extend, and which claims remain
conditional on them. That is the honest accounting at the door: not a passage
from finite work to continuum victory, but a careful list of what must be added
before continuum-facing statements can be read as more than finite shadows.

The bridge also keeps Chapter 16 from becoming a repetition of Chapter 15. It
does not need to reconstruct bounded atoms, redo compatibility, or recode finite
conditions. Those credentials have been filed. What remains is to test how far
they can be carried once the book permits completion language. The finite
apparatus gives Chapter 16 its skeleton; the next chapter must say which
continuum-facing assumptions put analytic flesh on that skeleton and which
claims still depend on those assumptions.

The chapter can therefore hand off without shame and without inflation. Behind
it are successor counts, bounded finite conditions, extension and compatibility,
finite decision, extensional coding, constructive countability, and the
mixed-pair residue tied back to the counting floor. Ahead of it are completion
hypotheses, continuum certificate families, boundary traces, and conditional
radiation. The bridge is the line between those two ledgers. It says: the finite
house is in order; the completion door is next; do not confuse one for the
other.

\section*{Coda: The Half-Filled Grid}

A constructor filling a crossword grid works the way this chapter's forcing
works: cell by cell, mark by mark, with a large picture in mind and only a
finite sheet of commitments actually on the table. The constructor does not
begin with the finished grid. She begins with a few letters, a few clues, a
promise not to write two letters in one square, and a willingness to extend the
partial work when a clue becomes decidable. Watching the grid fill is watching
the whole chapter in miniature: counting, bounded commitment, extension,
compatibility, merge, decision, coding, and the refusal to mistake finite
progress for the completed object.

A half-filled grid is not a failed grid. It is a legitimate state of the work.
Some cells carry letters; most are blank. The written cells are the finite
record. The blank cells are not errors or contradictions; they are the part of
the puzzle not yet committed. That is exactly what a finite spline condition was
in the chapter. It was not the spline, not a total assignment, and not a
continuum object. It was a finite list of bounded atom commitments: this knot,
this value, below the tolerance. The coda's grid gives the same idea in a form
the eye can hold. A condition is the set of squares already filled in pencil.

The first chapter task was to make sure the pencil could count its own marks.
Before the constructor can say that a grid has three committed cells, or that a
move adds one more square, she has to count. Chapter 15 therefore began at the
floor: successor, one, two, three, and the small certificate that the apparatus
can recognize its own finite successor steps. In the grid image, this is the
difference between a vague impression that some squares are filled and a ledger
that can say which finite marks have actually been made. The forcing calculus
does not float above arithmetic. It stands on the ability to count finite
commitments.

Filling one more cell extends the grid. The new partial grid contains every
letter the old partial grid contained, plus one more letter in one more square.
Nothing is erased. Nothing is reinterpreted. The new grid is stronger because it
has more information and still preserves the old information. That is extension
in its reader-facing form. A condition grows by adding bounded commitments while
keeping the commitments it already made. In the grid, this discipline is
obvious: a constructor who keeps erasing the old letters has not extended a
grid; she has changed the subject.

Two half-filled grids are compatible when they agree wherever both have written
something. If one grid writes \(A\) in a square and the other also writes \(A\)
there, no conflict appears. If one grid leaves that square blank, no conflict
appears either; the written letter can be carried into the merge. But if one
grid writes \(A\) and the other writes \(B\) in the same square, no single grid
can honor both. Compatibility is agreement on overlap. Merge is the finite act
of putting together two compatible records and keeping all their letters. The
least common stronger grid is the one that adds no extra letters beyond what
the two compatible partial grids already require.

This is why the chapter could treat decision as finite. A particular square can
be asked about. The constructor can point to that square and extend the current
partial grid by filling it, if a bounded admissible letter is being committed.
If she asks about finitely many squares, she can run finitely many such moves.
The result is a stronger partial grid that decides the finite question she
actually asked. The completed solution is not needed for that operation. Dense
decision in the chapter had the same modest force: every finite bounded
question the apparatus actually asks can be settled by finitely many bounded
extensions.

The partial grids also make extensional equality almost tactile. Two
constructors may have arrived at the same visible grid by different routes.
One solved the upper-left corner first; another solved the long central word
and then filled the corners. Their notebooks may contain different scratch
work. Their explanations may package the evidence differently. But if the same
cells carry the same letters, the partial grid is the same for the purpose of
the puzzle. That is extensional equality for proof-carrying conditions. The
certificate route does not become another letter in the square. What matters is
the finite record of actual commitments.

Once the grid is understood extensionally, it can be coded. A filled square is
finite data: a row, a column, and a letter. A half-filled grid is a finite list
of such data. Finite data can be placed in a finite code, and finite lists of
finite codes can themselves be coded by natural numbers. That is the chapter's
constructive countability in a picture. The family being counted is not the set
of completed puzzles. It is the family of finite partial fills. There are
countably many ways to write finitely many letters in finitely described cells,
even though the completed object toward which the work points belongs to a
different discussion.

The grid image must also remember the gauge side. A puzzle has local clues that
close cleanly and intersections where something remains to be reconciled. In
the chapter's gauge bookkeeping, the linear/vacuum part closes under the finite
certificate, while the mixed pair remains as the residue. The coda should not
turn that into a new analytic theorem. It is closer to saying: after the
constructor has honored all the finite entries licensed by the current ledger,
one crossing still carries the distinctive leftover mark. The mixed-pair
residue is filed because the finite bookkeeping can say what has closed and
what remains. It is not filed because the completed grid has already appeared.

This matters emotionally as well as formally. A half-filled grid can feel
unfinished because it is unfinished. But unfinished is not the same as
worthless. The chapter has spent many pages teaching the reader to respect the
power of finite work without flattering it. A finite condition is real. A
finite extension is progress. A compatibility test is information. A merge is a
new, stronger finite state. A decision of a finite question is an honest
settlement of that question. A natural code for a finite record is a genuine
counting result. The apparatus is allowed to be proud of all that, so long as it
does not confuse those achievements with the finished puzzle.

The completed grid is a different kind of object. It is the total fill in which
every square has its final letter and every clue is simultaneously satisfied.
The constructor may be guided by that object, may approach it through finite
work, and may speak about it as the goal. But none of the local moves already
possesses it. In the chapter's terms, the completed grid is the generic object
or completion-side object beyond the finite door. The finite forcing calculus
can approach it by extension and decide any finite bounded question on the way.
It cannot simply declare that the total object has been supplied.

The constructor's patience is the chapter's method. A good constructor does not
look at a half-filled grid and say that every blank square has secretly been
settled. She also does not despise the half-filled grid because it is not yet
complete. She keeps the two truths together. The marks already written are
binding. The blank squares remain open. The next admissible mark must respect
the old marks. If two drafts disagree at one square, they cannot be combined
until that conflict is resolved. This is the mood of the finite apparatus:
serious about what has been certified, cautious about what has not.

That is the finite/continuum boundary in its simplest image. A half-filled grid
is not a completed grid, even if every mark on it is correct. A countable list
of partial grids is not an enumeration of completed continuum objects. A
finite residue closeout is not an unconditional boundary-radiation theorem. A
separable pre-Hilbert skeleton is not a completed Hilbert space. A finite energy
ledger is not a Sobolev completion. The boundary is not a rebuke to the finite
apparatus. It is the condition under which its claims remain truthful.

The grid also explains why Chapter 16 must come next rather than merely
continue the coda. Once enough finite marks have been made, the natural question
changes. It is no longer only, ``Can this square be filled?'' It becomes,
``Under what hypotheses can the whole grid be treated as an object?'' That is a
different question. It asks about completion, limits, traces, and the structures
that let finite pieces speak as parts of a completed whole. Chapter 15 has
prepared the finite side of that question. It has not answered the completion
side.

So the coda closes Chapter 15 with a modest confidence. The apparatus has filled
many squares. It can count its marks. It can add new marks without erasing old
ones. It can test whether two partial fills agree. It can merge compatible
fills. It can decide any square it is actually asked to decide, and any finite
set of such squares by finite iteration. It can identify partial fills that
write the same letters, regardless of route. It can code the finite records and
list the finite possibilities. It can tie the gauge residue back to the finite
counting floor.

What it cannot do, and what this chapter refuses to pretend it can do, is hold
the completed grid merely because many finite pieces are in order. That work is
next. The next chapter must name the completion hypotheses, say what analytic
structure they add, and state which boundary-radiation claims remain
conditional on them. Chapter 15 leaves the reader at the table with a grid that
is genuinely fuller than before, governed by honest rules, and still not the
finished puzzle. The pencil has done real work. The door to completion is still
a door.

That is a better ending than a premature solution. The reader can see the
letters already placed. The reader can see how another letter would be placed.
The reader can see when two partial fills agree, when they cannot agree, and
why a finite list of requested squares can be settled without needing every
other square. The reader can also see the blank space that remains. Chapter 15
does not hide that blank space. It frames it, counts what is on this side of it,
and sends the book forward with the right question: what must be added before
the half-filled grid may be read as a completed one?
